% Copyright (c) 2026 Selina Cheggour. Tous droits réservés. Voir le fichier LICENSE.
% Parties 0 à 2 du notebook : mise en route, lire et contrôler le fichier, transformer en vecteurs.

\bandeaupartie{I. Parties 0 à 2 : lire, nettoyer, transformer en nombres}{Avant de regrouper quoi que ce soit, on vérifie ce qu'on a reçu.}

\section{Partie 0 : mise en route}

\begin{stepbox}[title={\faPlay\quad Trois cellules, rien à modifier}]
\begin{tabularx}{\linewidth}{>{\bfseries\color{sectioncolor}}c X}
0.1 & \faDatabase\ \texttt{NOM\_FICHIER} : le fichier de la base réelle, déjà écrit. \texttt{AFFICHER\_TEXTES} vaut \texttt{False} : aucun verbatim ne s'affiche. \\
0.2 & \faPuzzlePiece\ installe \texttt{py3langid}, le détecteur de langue, et \texttt{babel}, pour le nom des langues en clair. \\
0.3 & \faToolbox\ charge les outils : \texttt{verifier}, la palette des graphiques, \texttt{montrer} qui range chaque figure dans \texttt{figures/}. \\
\end{tabularx}
\doisvoir{\texttt{Base : RÉELLE (aucun verbatim ne sera affiché)} et le nom du fichier réel, puis \texttt{Outils prêts.}}
\end{stepbox}

\section{Partie 1 : lire le fichier sans le connaître à l'avance}

L'export de l'entreprise est un tableau Excel qui n'est pas \og propre \fg{}.

\begin{center}
\resizebox{\linewidth}{!}{%
\begin{tikzpicture}[x=1cm, y=0.58cm, font=\small]
\foreach \c/\x/\w in {A/0/2.2,B/2.2/6.4,C/8.6/1.9,D/10.5/2.2,E/12.7/1.6,F/14.3/1.2} {
  \node[font=\footnotesize\bfseries, color=slate] at (\x+\w/2, 0.6) {\c};
}
\foreach \r in {1,...,9} \node[font=\footnotesize\bfseries, color=slate] at (-0.4,-\r+0.5) {\r};
\fill[notebg] (0,-2) rectangle (15.5,0); \node[color=slate, font=\footnotesize\itshape] at (7.7,-1) {lignes 1 et 2 : vides};
\fill[sectioncolor!25] (0,-3) rectangle (15.5,-2);
\foreach \t/\x/\w in {Case Number/0/2.2,Verbatim/2.2/6.4,Brand/8.6/1.9,Business Segment/10.5/2.2,Year Period/12.7/1.6,Blank/14.3/1.2}
  \node[font=\footnotesize\bfseries, color=darkblue] at (\x+\w/2,-2.5) {\t};
\foreach \r/\num/\verb/\mq/\seg/\per in {3/10001234/{EXTERNALBonjour,j'ai acheté...}/Brand 003/Chocolate/26 P05,
  4/10001302/{EXTERNALHello,the bar *******...}/Brand 007/Chocolate/26 P05,
  5/10001377/{EXTERNALGuten Tag,die Packung...}/Brand 012/Food/26 P06} {
  \node[anchor=west, font=\scriptsize\ttfamily] at (0.05,-\r-0.5) {\num};
  \node[anchor=west, font=\scriptsize\ttfamily] at (2.25,-\r-0.5) {\verb};
  \node[anchor=west, font=\scriptsize] at (8.65,-\r-0.5) {\mq};
  \node[anchor=west, font=\scriptsize] at (10.55,-\r-0.5) {\seg};
  \node[anchor=west, font=\scriptsize] at (12.75,-\r-0.5) {\per};
  \node[font=\scriptsize] at (14.9,-\r-0.5) {$\blacksquare$};
}
\node[color=slate, font=\footnotesize\itshape] at (7.7,-6.5) {\ldots\ des centaines de lignes \ldots};
\fill[notebg] (0,-9) rectangle (15.5,-7); \node[color=slate, font=\footnotesize\itshape] at (7.7,-8) {lignes vides en fin de tableau};
\draw[slate!50] (0,-9) grid[xstep=15.5, ystep=1] (15.5,0);
\foreach \x in {2.2,8.6,10.5,12.7,14.3} \draw[slate!50] (\x,0) -- (\x,-9);
\draw[accentblue, very thick, rounded corners] (14.25,-9.1) rectangle (15.55,0.1);
\node[accentblue, font=\footnotesize\bfseries, anchor=north] at (14.9,-9.2) {à ignorer};
\draw[-{Stealth}, accentblue, thick] (-1.6,-2.5) node[left, font=\footnotesize\bfseries, text width=1.6cm, align=right] {l'en-tête est en ligne 3} -- (-0.7,-2.5);
\end{tikzpicture}}
\end{center}
{\small\color{slate}\textit{Schéma de la disposition du fichier ; les lignes montrées sont inventées.}}

\begin{stepbox}[title={1.1\quad Un lecteur qui trouve tout seul où commencent et finissent les données}]
\texttt{charger\_plaintes(chemin)} lit le fichier \textbf{sans en-tête et en texte}, cherche la première
ligne qui contient à la fois \texttt{Case Number} et \texttt{Verbatim}, garde les cinq colonnes utiles,
ignore \texttt{Blank} et les colonnes sans nom, retire les lignes vides, et rend un petit
\textbf{rapport} (feuille, ligne d'en-tête, lignes vides retirées, colonnes ignorées).
\doisvoir{l'en-tête trouvé tout seul (en \textbf{ligne 3} sur l'export), le nombre de plaintes lues, la colonne \texttt{Blank} ignorée, les cinq colonnes attendues et seulement elles.}
\image{un facteur qui cherche le nom sur l'enveloppe, au lieu de compter \og troisième enveloppe en partant du haut \fg.}
\end{stepbox}

\begin{piege}
\texttt{pd.read\_excel(chemin, skiprows=2, nrows=1301)} marche aujourd'hui, et \textbf{cassera en silence}
demain : une ligne de titre en plus ou un fichier plus long, et on lit de travers sans aucun message.
\end{piege}

\reflexion{1.1}{Si la prochaine livraison compte 40\,000 lignes, avec une ligne de titre en plus avant
l'en-tête, ou arrive en CSV avec un séparateur \texttt{;}, que faut-il changer dans
\texttt{charger\_plaintes} ? Et dans la version \texttt{skiprows=2, nrows=1301} ?}

\begin{stepbox}[title={1.2\quad Nettoyer les verbatims, dans le bon ordre}]
\begin{center}
\begin{tikzpicture}[
  t/.style={draw=#1, fill=#1!8, rounded corners=4pt, text width=13.8cm, align=left, font=\footnotesize\ttfamily, inner sep=5pt},
  e/.style={font=\footnotesize, color=slate, text width=13.8cm, align=center}]
\node[t=slate] (a) {EXTERNALBonjour,j'ai acheté une tablette ******* hier.Elle avait un goût de savon !Lien : https://exemple.org/fiche.produit};
\node[e, below=0.35cm of a] (s) {1. retirer \texttt{EXTERNAL} \quad 2. retirer le lien \quad 3. retirer les étoiles\\ 4. une espace après la ponctuation \quad 5. des espaces simples};
\node[t=accentblue, below=0.35cm of s] (b) {Bonjour, j'ai acheté une tablette hier. Elle avait un goût de savon ! Lien :};
\draw[-{Stealth[length=3mm]}, very thick, color=accentblue] (s.west) ++(-0.2,0.3) -- ++(0,-0.6);
\end{tikzpicture}
\end{center}
{\small\color{slate}\textit{Exemple inventé, dans la forme des verbatims de la base fictive.}}
\doisvoir{plus aucun texte ne commence par \texttt{EXTERNAL}, plus d'étoiles ni de liens, les phrases sont séparées, une adresse web n'est pas coupée en morceaux.}
\image{on enlève l'étiquette collée sur la lettre avant de la lire, et on sépare les mots écrits sans espace.}
\end{stepbox}

\reflexion{1.2}{Pourquoi retirer le lien \textbf{avant} de remettre des espaces après la ponctuation ?}

\begin{stepbox}[title={1.3\quad Contrôler ce qu'on a reçu}]
Doublons de \texttt{Case Number}, verbatims vides, longueur des textes en mots, comptes par marque, segment et période.
\doisvoir{aucun doublon, aucun verbatim vide ; la longueur médiane des plaintes, puis le nombre de marques, de segments et de périodes.}
\end{stepbox}

\illus[0.97\linewidth]{1_plaintes_par_periode.png}{Figure 1 du notebook (base fictive) : le nombre de plaintes par période.}

\begin{lire}
\textbf{Comment lire ce graphique.} Une barre par période (\og 26 P05 \fg{} : année 2026, cinquième période).
La ligne pointillée est la moyenne ; les barres rouges la dépassent de plus d'un quart. Ce graphique dit
\textbf{quand} il y a eu beaucoup de plaintes, jamais \textbf{pourquoi}.
\end{lire}

\reflexion{1.3}{Le graphique montre des bosses, l'été 2025 et au printemps 2026. Est-ce un incident ?}

\begin{stepbox}[title={1.4\quad Quelles langues ?}]
\texttt{py3langid.classify(texte)} rend le code de la langue et un score. La base n'a ni colonne pays ni
colonne langue : la langue est le seul indice de géographie.
\doisvoir{une langue pour chaque plainte, et \textbf{plusieurs langues} (les plaintes viennent d'Europe et du Moyen-Orient), de la plus fréquente à la plus rare, avec leur nom en clair. Quelques langues très rares sont souvent des erreurs du détecteur sur des textes courts.}
\image{la langue d'une lettre dit dans quelle boîte de tri la mettre, pas de quelle ville elle vient : l'allemand couvre plusieurs pays.}
\end{stepbox}

\begin{stepbox}[title={1.5\quad L'anonymisation a-t-elle laissé passer quelque chose ?}]
On \textbf{compte}, sans lire les textes, les verbatims d'origine qui contiennent encore une forme reconnaissable :
\begin{center}\small
\renewcommand{\arraystretch}{1.25}
\begin{tabular}{ll}
\toprule
\textbf{motif} & \textbf{expression régulière} \\
\midrule
adresse électronique & \verb|[\w.+-]+@[\w-]+\.[\w.]+| \\
numéro de téléphone & \verb|\+?\d[\d .-]{7,}\d| \\
suite d'au moins six chiffres & \verb|\b\d{6,}\b| \\
adresse web & \verb|https?://\S+| \\
\bottomrule
\end{tabular}
\end{center}
\doisvoir{ni courriel, ni téléphone, ni longue suite de chiffres ; \textbf{40 plaintes} contiennent un lien web ; toutes citent un produit masqué par des étoiles.}
\end{stepbox}

\reflexion{1.5}{Sur la base réelle, un compte non nul ne prouve pas une fuite, et un compte nul ne prouve
pas l'absence de fuite. Pourquoi ?}

%%%%%%%%%%%%%%%%%%%%%%%%%%%%%%%%%%%%%%%%%%%%%%%%%%%%%%%%%%%%%%%%%%%%%%%%%%%%%%%%%%%%%%%%%%%%%%%%%%%%%%%%
\clearpage
\section{Partie 2 : transformer chaque plainte en vecteur}

Pour regrouper des textes, il faut les transformer en nombres. Deux familles s'opposent.

\begin{center}
\resizebox{\linewidth}{!}{%
\begin{tikzpicture}[font=\small]
\begin{scope}
\node[font=\bfseries, color=sectioncolor] at (2.2,4.4) {TF-IDF : un axe par mot};
\draw[-{Stealth}, thick, slate] (0,0) -- (4.2,0) node[right] {\texttt{savon}};
\draw[-{Stealth}, thick, slate] (0,0) -- (0,3.8) node[above] {\texttt{Seife}};
\draw[-{Stealth}, thick, slate] (0,0) -- (-1.9,-1.6) node[below left] {\texttt{soap}};
\fill[accentblue] (3.2,0) circle (3pt) node[above, font=\footnotesize] {plainte FR};
\fill[ambre] (0,2.9) circle (3pt) node[right, font=\footnotesize] {plainte DE};
\fill[darkblue] (-1.4,-1.18) circle (3pt) node[right, font=\footnotesize] {plainte EN};
\node[text width=5.4cm, align=center, font=\footnotesize, color=slate] at (2.0,-2.7) {trois mots sans rapport : les trois plaintes sont à angle droit, aussi loin que possible};
\end{scope}
\begin{scope}[xshift=9.3cm]
\node[font=\bfseries, color=accentblue] at (1.8,4.4) {plongement : un espace de sens};
\fill[notebg] (1.8,1.7) ellipse (2.3 and 1.6);
\node[font=\footnotesize\itshape, color=slate] at (1.8,3.0) {\og goût de savon \fg};
\fill[accentblue] (1.4,1.6) circle (3pt) node[left, font=\footnotesize] {FR};
\fill[ambre] (2.1,1.9) circle (3pt) node[right, font=\footnotesize] {DE};
\fill[darkblue] (1.8,1.2) circle (3pt) node[below, font=\footnotesize] {EN};
\fill[slate] (4.4,-0.6) circle (3pt) node[right, font=\footnotesize] {\og emballage abîmé \fg};
\node[text width=5.4cm, align=center, font=\footnotesize, color=slate] at (1.9,-2.7) {un modèle entraîné sur plus de cent langues place les trois plaintes côte à côte};
\end{scope}
\end{tikzpicture}}
\end{center}

\begin{stepbox}[title={2.1\quad TF-IDF, et à quoi ressemblent les plus proches voisins}]
Chaque plainte $a$ devient un vecteur avec une coordonnée par terme $t$ :
\begin{formulabox}
\[
w_{t,a} = \bigl(1 + \ln \mathrm{tf}_{t,a}\bigr)\left(\ln\frac{1+n}{1+n_t} + 1\right)
\quad\text{si } \mathrm{tf}_{t,a} > 0, \qquad w_{t,a} = 0 \text{ sinon},
\]
où $\mathrm{tf}_{t,a}$ est le nombre de fois où $t$ apparaît dans la plainte $a$, $n$ le nombre de plaintes et
$n_t$ le nombre de plaintes qui contiennent $t$. Le vecteur est ensuite divisé par sa longueur. La
proximité de deux plaintes est le \textbf{cosinus} de leurs vecteurs $\mathbf{x}_a$, $\mathbf{x}_b$ :
$\cos(\mathbf{x}_a,\mathbf{x}_b) = \mathbf{x}_a^{\top}\mathbf{x}_b / (\lVert\mathbf{x}_a\rVert\,\lVert\mathbf{x}_b\rVert)$,
qui vaut 1 pour deux textes identiques et 0 pour deux textes sans mot commun.
\end{formulabox}
\illus[0.62\linewidth]{ill_tfidf.png}{Exemple inventé : un mot rare et répété pèse lourd, un mot partout (\og le \fg) presque rien.}
Ensuite, pour chaque plainte, on cherche sa plus proche voisine et on compte la part des plaintes dont la
voisine est \textbf{dans la même langue}.
\doisvoir{\textbf{757 mots} dans le vocabulaire ; voisine dans la même langue : \textbf{100 \%}.}
\end{stepbox}

\begin{stepbox}[title={2.2\quad Plongements multilingues : un transformer qui lit la plainte}]
\begin{center}
\resizebox{\linewidth}{!}{%
\begin{tikzpicture}[
  bx/.style={draw=#1, fill=#1!10, rounded corners=4pt, align=center, font=\footnotesize, minimum height=1.3cm},
  fl/.style={-{Stealth[length=2.5mm]}, thick, color=slate}]
\node[bx=slate, text width=3.2cm] (t) {\og Elle avait un goût de savon ! \fg};
\node[bx=sectioncolor, right=0.5cm of t, text width=2.8cm] (j) {jetons\\\texttt{Elle|avait|un|}\\\texttt{goût|de|sav|on|!}};
\node[bx=darkblue, right=0.5cm of j, text width=2.6cm] (m) {{\Large\faLayerGroup}\\28 couches d'attention};
\node[bx=accentblue, right=0.5cm of m, text width=2.6cm] (v) {un vecteur de\\\textbf{1\,024 nombres}};
\draw[fl] (t) -- (j); \draw[fl] (j) -- (m); \draw[fl] (m) -- (v);
\end{tikzpicture}}
\end{center}
Le modèle \texttt{Qwen/Qwen3-Embedding-0.6B} (Apache 2.0, plus de cent langues) reçoit chaque plainte
précédée d'une consigne (\og identifie le problème décrit \fg) et rend un vecteur de longueur 1. Dans
chaque couche, chaque jeton \og regarde \fg{} tous les autres pour préciser son sens : c'est
l'\textbf{attention}. Si le modèle ne se charge pas, la cellule passe à \texttt{multilingual-e5-base}.
\doisvoir{un vecteur normé par plainte ; voisine dans la même langue : \textbf{100 \%}, comme le TF-IDF.}
\image{un interprète qui résume chaque lettre en une \og empreinte de sens \fg, quelle que soit sa langue.}
\end{stepbox}

\reflexion{2.2}{Le plongement multilingue rapproche-t-il vraiment des plaintes de langues différentes ?
Regarde le pourcentage.}

\begin{stepbox}[title={2.3\quad Retirer la signature de la langue : le centrage}]
\begin{formulabox}
\[
\tilde{\mathbf{e}}_i = \frac{\mathbf{e}_i - \boldsymbol{\mu}_{\ell(i)}}{\lVert \mathbf{e}_i - \boldsymbol{\mu}_{\ell(i)} \rVert},
\]
où $\mathbf{e}_i$ est le plongement de la plainte $i$, $\ell(i)$ sa langue, et $\boldsymbol{\mu}_\ell$ la moyenne
des plongements des plaintes de la langue $\ell$. Une langue de moins de 10 plaintes reçoit la moyenne de
toutes les plaintes.
\end{formulabox}
\illus[0.92\linewidth]{ill_centrage.png}{Exemple inventé : la couleur est la langue, la forme le sujet. À gauche, on voit d'abord les langues ; à droite, après centrage, les sujets.}
\image{trois personnes qui parlent fort, doucement et avec un accent : on retire à chacune sa façon de
parler, et il ne reste que ce qu'elle dit.}
\doisvoir{un vecteur centré et normé par plainte ; voisine dans la même langue : \textbf{98 \%}, un peu moins qu'avant.}
\end{stepbox}

\illus[0.98\linewidth]{2_carte_langues_avant_apres.png}{Figure 2 du notebook (base fictive) : les plaintes en deux dimensions (t-SNE), avant et après centrage.}

\begin{lire}
\textbf{Comment lire ce graphique.} Chaque point est une plainte, la couleur est sa langue. Le t-SNE
place près l'une de l'autre deux plaintes proches dans l'espace des vecteurs ; \textbf{les axes n'ont pas
de sens}, seules les distances en ont. À gauche, les couleurs forment des îlots séparés. À droite, si le
centrage marche, les trois couleurs se mélangent dans les mêmes paquets : chaque paquet est un sujet
écrit dans plusieurs langues.
\end{lire}

\reflexion{2.3}{Le centrage retire-t-il seulement la langue ?}
